% Articulación posterior a la generación de pi, phi y e.
% Contenido acumulativo; no reemplaza el desarrollo previo del TRIT.
\subsection{Realizaciones exponenciales de las coordenadas generadas}
\label{euler-trit-articulacion}

Las coordenadas regionales ya determinadas permiten reconocer las
realizaciones concretas de la colección cuadrática del TRIT. El orden es
sustantivo: la construcción de las regiones precede a la evaluación de una
exponencial en los parámetros \(\pi\) o \(2\log\varphi\). Estas expresiones
identifican la función de las coordenadas en operadores construidos desde
APP y TPK; no intervienen en la selección de los prefijos que las producen.

\subsubsection{Ciclo ternario, transporte nilpotente e incidencia de hojas}

La multiplicación \(a(r)=4r\) en \(\mathbb Z/9\mathbb Z\) tiene orden
\(3\). Sobre las unidades determina las órbitas \(\{1,4,7\}\) y
\(\{2,8,5\}\). En el módulo de aumento de cualquiera de ellas, formado por
las combinaciones enteras cuyos coeficientes suman cero, la base
\((e_r-e_{4r},e_{4r}-e_{7r})\) da
\[
 C=\begin{pmatrix}0&-1\\1&-1\end{pmatrix},
 \qquad C^2+C+I=0.
\]
La forma restringida tiene matriz de Gram
\(\bigl(\begin{smallmatrix}2&-1\\-1&2\end{smallmatrix}\bigr)\).
En las evaluaciones residuales módulo \(9\), la acción distingue las hojas:
\(S(ax,ay)=aS(x,y)\), mientras \(P(ax,ay)=a^{-1}P(x,y)\).
La realización cíclica conserva, por tanto, una orientación contragrediente
entre suma y producto. Los cocientes del levantamiento entero se conservan
en el registro aritmético correspondiente.

El transporte parabólico elemental se representa mediante
\[
 N=\begin{pmatrix}0&1\\0&0\end{pmatrix}.
\]
La incidencia areal--volumétrica ya construida en \eqref{eq:fav} proporciona
\[
 F_{\rm av}=\begin{pmatrix}0&1\\1&1\end{pmatrix},
 \qquad A=F_{\rm av}^2=\begin{pmatrix}1&1\\1&2\end{pmatrix}.
\]
Los tres operadores actúan en espacios reales declarados: el módulo de
aumento tensorizado con \(\mathbb R\), el plano del transporte nilpotente y
el plano de incidencia modal. Compartir dimensión matricial no identifica
estos espacios ni sus coordenadas.

\begin{proposition}[Generadores concretos de los tres tipos]
\label{euler-trit-generadores}
Los operadores
\[
 J_{\mathrm e}=\frac{2C+I}{\sqrt3},\qquad
 J_{\mathrm p}=N,\qquad
 J_{\mathrm h}=\frac{2A-3I}{\sqrt5}
\]
satisfacen, respectivamente,
\[
 J_{\mathrm e}^{\,2}=-I,\qquad
 J_{\mathrm p}^{\,2}=0,\qquad
 J_{\mathrm h}^{\,2}=I.
\]
\end{proposition}
\begin{proof}
La relación de \(C\) implica \((2C+I)^2=-3I\).
La multiplicación directa da \(N^2=0\).
Finalmente, \(A\) tiene traza \(3\) y determinante \(1\), por lo que
\(A^2-3A+I=0\) y \((2A-3I)^2=5I\).
\end{proof}

\begin{proposition}[Tres evaluaciones de la identidad exponencial]
\label{euler-trit-evaluaciones}
En las realizaciones anteriores,
\begin{equation}
 \boxed{
 \exp(\pi J_{\mathrm e})+I=0,\qquad
 \exp(J_{\mathrm p})=I+J_{\mathrm p},\qquad
 \exp(2\log\varphi\,J_{\mathrm h})=F_{\rm av}^{\,2}.}
 \label{euler-trit-tres-evaluaciones}
\end{equation}
\end{proposition}
\begin{proof}
La especialización elíptica de la exponencial trítica da
\(\cos\pi\,I+\sin\pi\,J_{\mathrm e}=-I\).
La nilpotencia elimina exactamente los términos de grado al menos \(2\)
en la segunda expresión. Para la tercera, \(\varphi^2=\varphi+1\) implica
\[
 \cosh(2\log\varphi)=\frac32,\qquad
 \sinh(2\log\varphi)=\frac{\sqrt5}{2}.
\]
Por tanto, la exponencial vale
\(\frac32I+\frac{\sqrt5}{2}J_{\mathrm h}=A\).
\end{proof}

La primera igualdad representa la media vuelta elíptica; la vuelta completa
satisface \(\exp(2\pi J_{\mathrm e})=I\). La segunda mantiene un término
lineal no nulo: el régimen parabólico expresa transporte, no ausencia de
evolución. La tercera realiza un avance hiperbólico unimodular. Sus
autovalores son \(\varphi^2\) y \(\varphi^{-2}\); una dirección se expande y
la otra se contrae conservando el determinante. Las tres evaluaciones
emplean parámetros distintos de una colección exponencial común.

En esa colección, \(e\) interviene mediante la evaluación escalar
\(\exp(I)=eI\) y la ley \(\exp(sI)\exp(tI)=\exp((s+t)I)\).
Su papel no se restringe al tipo parabólico. La coordenada \(e\), obtenida
previamente de su región, es reconocida aquí como la evaluación unitaria
de la exponencial; esta identidad no sustituye su genealogía coinductiva.

\subsubsection{Resolución polinómica de los regímenes}

Las tres realizaciones admiten una presentación compacta sobre la suma
directa de sus espacios. Defínanse
\[
 \mathbb J=J_{\mathrm e}\oplus J_{\mathrm p}\oplus J_{\mathrm h},
 \qquad Q=\mathbb J^2.
\]
Las relaciones anteriores implican
\[
 \mathbb J^6=\mathbb J^2,\qquad Q^3=Q.
\]
El uso de \(Q\) mantiene separado este cuadrado operatorial del registro
dodecafásico \(K\).

\begin{proposition}[Proyectores polinómicos de régimen]
\label{euler-trit-proyectores}
Las aplicaciones
\[
 p_{\mathrm e}=\frac{Q^2-Q}{2},\qquad
 p_{\mathrm p}=I-Q^2,\qquad
 p_{\mathrm h}=\frac{Q^2+Q}{2}
\]
son idempotentes mutuamente ortogonales, suman \(I\) y recuperan los tres
sumandos de la representación.
\end{proposition}
\begin{proof}
Los tres polinomios toman los valores \((1,0,0)\), \((0,1,0)\) y
\((0,0,1)\) en los puntos \(-1,0,+1\), respectivamente. Las diferencias
entre sus cuadrados y ellos mismos, así como sus productos cruzados, son
múltiplos de \(X^3-X\). Sustituir \(Q\), que anula ese polinomio, demuestra
las identidades. En los tres bloques, \(Q\) actúa como \(-I,0,I\).
\end{proof}

Se obtiene así
\begin{align}
 \exp(t\mathbb J)
 ={}&p_{\mathrm e}(\cos t\,I+\sin t\,\mathbb J)
   +p_{\mathrm p}(I+t\mathbb J)\notag\\
   &+p_{\mathrm h}(\cosh t\,I+\sinh t\,\mathbb J).
 \label{euler-trit-exponencial-total}
\end{align}
Cada sumando conserva su relación cuadrática y la conjugación invierte
su evolución. La representación total reúne los tres regímenes, pero no
define por sí sola un operador de transición entre ellos. Tal transición
requiere los mapas de transporte del TPK con sus dominios, orientación y
memoria. La articulación exponencial permite reconocer sus tipos sin
reemplazar esa dinámica por una identificación de las fibras.
